% Krawatzek & Pestel: New Histories of and for Europe: Reading Summary
% Introduction to the European Union, Week 2
% Compile with: pdflatex krawatzek-pestel-new-histories-europe-reading-summary.tex

\documentclass[11pt,a4paper]{article}
\usepackage[utf8]{inputenc}
\usepackage[T1]{fontenc}
\usepackage[margin=2.5cm]{geometry}
\usepackage{booktabs}
\usepackage{enumitem}
\setlist{nosep,leftmargin=*}
\usepackage{parskip}
\usepackage{microtype}
\usepackage{hyperref}
\usepackage{xcolor}
\definecolor{eublue}{RGB}{0,51,153}

\hyphenation{Spitzenkandidat Spitzenkandidaten}
\hypersetup{
  colorlinks=true,
  linkcolor=eublue,
  urlcolor=eublue,
  pdftitle={Krawatzek \& Pestel: New Histories of Europe},
  pdfauthor={Introduction to the EU, UCM}
}

\begin{document}

\title{Krawatzek \& Pestel: New Histories of and for Europe\\Narrating the European Project}
\author{Introduction to the European Union\\Universidad Complutense de Madrid\\Week 2\\\textit{Contemporary European History} (2023)}
\date{}
\maketitle
\vspace{-1em}
\hrule
\vspace{1em}

\section*{Rhetorical Framework (Ethos, Logos, Pathos)}

\begin{table}[h]
\centering
\begin{tabular}{@{}lll@{}}
\toprule
\textbf{Appeal} & \textbf{Meaning} & \textbf{Krawatzek \& Pestel's Use} \\
\midrule
\textbf{Ethos} & Credibility, authority, trust & Scholarly review; meta-critique of historiography; balanced assessment \\
\textbf{Logos} & Logic, structure, argument & Systematic comparison of approaches; identification of gaps and biases \\
\textbf{Pathos} & Emotion, stakes, persuasion & Implicit appeal for diversity, inclusion of marginal voices \\
\bottomrule
\end{tabular}
\end{table}

\textbf{Key insight:} The article uses \textit{logos} (categorizing and critiquing narratives) and \textit{ethos} (positioning as fair reviewers). The \textit{pathos} is indirect---advocacy for ``diverse and vernacular voices'' appeals to the reader's sense of fairness and inclusion.

\section*{Bibliographic Information}

\begin{itemize}
\item \textbf{Authors:} Felix Krawatzek \& Friedemann Pestel
\item \textbf{Title:} New Histories of and for Europe: Narrating the European Project
\item \textbf{Journal:} Contemporary European History
\item \textbf{Year:} 2023
\item \textbf{DOI:} 10.1017/S0960777323000450
\end{itemize}

\section*{Summary \& Key Points}

\subsection*{I. The Evolving European Project}

\begin{itemize}
\item The European project is constantly evolving, shaped by historical junctures and ongoing crises (WWI, WWII, Treaty of Rome, Brexit, war in Ukraine).
\item Recent scholarship aims to reach broader audiences, often making political interventions and judgments about Europe's past and present.
\item The European project is narrated through multiple lenses: intellectual, institutional, and peace-oriented perspectives.
\end{itemize}

\subsection*{II. Reviewed Scholarship (Rhetorical Profiles)}

\textbf{II.A. Pagden's Intellectual History} --- Traces Europe's roots from Enlightenment to postwar integration. Strong \textit{logos} (ideas as causes); \textit{ethos} (philosophical lineage); \textit{pathos} (identity, belonging).

\textbf{II.B. Ghervas's Peace-Oriented Approach} --- Centers on Europe's quest for peace. Strong \textit{pathos} (peace, fear of war); \textit{logos} (peace as causal driver); \textit{ethos} (moral purpose of integration).

\textbf{II.C. Gilbert's Institutional History} --- Focuses on institutional and political history. Strong \textit{logos} (institutions, structures); \textit{ethos} (elites, leaders); weaker \textit{pathos} (ordinary people absent).

\textbf{II.D. Van Middelaar's Insider View} --- Insider's view of EU institutions. Strong \textit{ethos} (insider authority); \textit{logos} (institutional mechanics); \textit{pathos} (legitimacy as emotional stake).

\textbf{II.E. Jarausch's Defensive Account} --- Defends achievements of European integration. Strong \textit{pathos} (defense, pride, contrast with US); \textit{ethos} (scholarly defense); \textit{logos} (comparative argument).

\subsection*{III. Critical Observations}

\begin{itemize}
\item All reviewed works tend to affirm the European project and its institutions, but often privilege top-down, pro-European perspectives.
\item The reviewed literature often overlooks more diverse or critical voices, including anti-European, Eurosceptic, and alternative perspectives.
\item Future scholarship should include more diverse and vernacular voices.
\end{itemize}

\section*{Main Arguments}

\begin{enumerate}
\item \textbf{Multiple Narratives} --- Europe narrated through intellectual, institutional, peace-oriented lenses.
\item \textbf{Top-Down Bias} --- Elites and institutions privileged; ordinary people's voices absent.
\item \textbf{Pro-European Bias} --- Critical and Eurosceptic perspectives underrepresented.
\item \textbf{Political Interventions} --- Scholarship not neutral; scholars as advocates.
\item \textbf{Constant Evolution} --- Crises reshape narratives.
\end{enumerate}

\section*{Context \& Analysis}

\textbf{Meta-note:} The \textit{reviewed} scholars (Pagden, Ghervas, Gilbert, etc.) each deploy ethos, logos, and pathos differently---K\&P analyze how those choices shape what gets told and what gets left out.

\section*{Relevance to Course}

\begin{itemize}
\item \textbf{Ethos:} Importance of critical perspectives; questioning official narratives.
\item \textbf{Logos:} Multiple ways to narrate European history; different causal stories.
\item \textbf{Pathos:} Value of diverse voices; stakes of whose story gets told.
\item \textbf{All three:} Critical engagement with sources---credibility, reasoning, and representation together.
\end{itemize}

\section*{Discussion Questions}

\begin{itemize}
\item How do different approaches shape understanding? (Ethos: whose authority? Logos: causal stories? Pathos: values emphasized?)
\item Why do top-down perspectives dominate? (Ethos: who has access? Logos: structural incentives? Pathos: whose voices matter?)
\item What perspectives are missing? (Ethos: who is excluded? Logos: absent evidence? Pathos: whose story untold?)
\item Should Eurosceptic/critical voices be included? (Ethos: fairness, balance; Logos: completeness; Pathos: representation)
\end{itemize}

\vspace{1em}
\noindent\footnotesize\textit{Reference:} Krawatzek, F., \& Pestel, F. (2023). New Histories of and for Europe: Narrating the European Project. \textit{Contemporary European History}, 1--10. DOI: 10.1017/S0960777323000450

\end{document}
